% latexmk -pdf -pvc blatt9.tex
\documentclass{hott-übung}

\begin{document}

\setcounter{blattnummer}{9}
\setcounter{aufgabennummer}{0}

\blatt

\aufgabe{}
Zeige: Ein punktierter Typ \(A\) ist genau dann zusammenhängend, wenn \(\|A\|_0\) kontrahierbar ist.


\aufgabe{}
Für \(n\geq -1\) ist \(|\_|_n:A\to \|A\|_n\) surjektiv.

\aufgabe{}
Sei \(A : \mU\) und \(B : \|A\|_n \to \mU\). Zeige:
\[\left\|\sum_{x:A}B(|x|_{n})\right\|_n\simeq \sum_{y:\|A\|_n}\|B(x)\|_n\]
\emph{Hinweis: Vergleiche mit dem Beweis von \(\|A \times B\|_n \simeq \|A\|_n \times \|B\|_n\).}

\aufgabe{}
Sei \((f,p_f) : (A,a) \to_\ast (B,b)\) eine punktierte Abbildung, also \(p_f : f(a) = b\).
Wir definieren
\begin{align*}
  \Omega f & : \Omega A \to \Omega B \\
  (\Omega f)(q) & \colonequiv p_f^{-1} \kon \ap_f(q) \kon p_f.
\end{align*}
% Außerdem sei \(p_{\Omega f} : p_f^{-1}\kon p_f = \refl_b\) gegeben (wie etwa in Lemma 1.5.10), sodass \((\Omega f,p_{\Omega f}) : \Omega A \to_\ast \Omega B\).
Zeige: Sind \(A\) und \(B\) zusammenhängend, so ist \(f\) genau dann eine Äquivalenz, wenn \(\Omega f\) eine ist.
\emph{Hinweis: Lemma 2.4.12.}

\aufgabe{}
Irgendwas mit \(S^n\)


\end{document}
